\chapter{Literature Review / Related Work}
This chapter reviews the research that motivates and enables FaceCon. Because FaceCon is a research and development project, a detailed and critical literature survey is provided rather than a catalogue of competing products. The review is organized topic by topic, moving from the biology of facial genetics through to the methods the system depends on. Within each topic, every research item is reviewed separately under three headings: a Summary of what the work proposes and finds, a Critical Analysis of its strengths and limitations, and a statement of its Relationship to the Proposed Work. Each topic then closes with a short synthesis that draws its papers together and connects them to FaceCon.

\section{Definitions, Acronyms, and Abbreviations}
This chapter uses the terms defined in Chapter~1, together with the following. Genotype is the set of genetic variants carried by an individual. Phenotype is an observable trait such as facial shape. Locus is a position in the genome. Heritability is the proportion of trait variation attributable to genetic variation. Mesh is a set of vertices and faces describing a 3D surface. Morphable model is a parametric model that represents a family of shapes as a low-dimensional vector. Contrastive learning is training that pulls matched pairs together and pushes unmatched pairs apart in a shared embedding space.

\section{Detailed Literature Review}
The review is grouped into topics that trace the path from the biology of facial genetics to the methods FaceCon depends on. Each topic opens with its most relevant works, reviews every item under the same three headings, and closes with a short synthesis that draws out the implications for FaceCon.

\subsection{Genetic Prediction of Facial Morphology}
This topic covers the genome-wide association studies that establish, biologically, that facial shape is partly heritable, and that identify the variants a predictive model can draw on.

\subsubsection{Claes et al. --- Global to Local Genetic Effects on Facial Shape}
Summary. Claes et al.~\cite{ref:claes} propose a data-driven, global to local approach for discovering how genetic variants influence facial shape. The objective is to move beyond simple landmark distances and to localize genetic effects to specific facial regions. Methodologically, the work hierarchically segments three-dimensional facial shape into progressively finer regions and performs a multivariate genome-wide association study on each segment, relating genome-wide SNPs to the shape of each region. The data combine three-dimensional facial images with genome-wide genotyping for cohorts of predominantly European ancestry. The study identifies multiple genomic loci associated with facial shape, several of them novel, and shows that distinct loci act on distinct parts of the face.

Critical Analysis. The principal strength of the work is methodological: the multivariate, region design captures effects that univariate landmark studies miss, and the resulting loci are published and reproducible. Its main limitation for a prediction project is that it is oriented toward discovery rather than generation, so it identifies associations without delivering an end-to-end face predictor. The cohort is ancestry-concentrated, which limits the transfer of specific effect sizes to other populations, and, as with all facial genetics, the effects are individually small and highly polygenic.

Relationship to the Proposed Work. FaceCon adopts the face-associated loci reported in this line of work as part of the evidence base for its documented SNP panel, rather than copying an undisclosed panel. It differs by building a generative genotype to FLAME predictor instead of an association map, and it directly inherits the study's caution about ancestry generalization, which it addresses through stratified evaluation. The work therefore influences FaceCon's SNP-selection stage rather than its modelling architecture.

\subsubsection{White et al. --- Genetic Architecture of the Human Face}
Summary. White et al.~\cite{ref:white} extend segment facial genetics with a multivariate genome-wide association study that expands the catalogue of face-associated variants and clarifies how genetic effects are distributed across the face. The approach again pairs three-dimensional facial phenotyping with genome-wide genotypes in predominantly European cohorts, and it contributes a richer map of which genomic regions influence which aspects of facial morphology.

Critical Analysis. The study's strength is the breadth of its catalogue and the rigor of its multivariate analysis, which strengthen the evidence that facial shape is genetically structured. Its limitations mirror those of the topic as a whole: it remains discovery-oriented rather than predictive, its cohorts are ancestry-concentrated, and the polygenic, small-effect nature of the signal means that individual-level prediction is not demonstrated.

Relationship to the Proposed Work. FaceCon uses this expanded catalogue as a second pillar of its SNP-panel evidence base, alongside Claes et al. It modifies the role of these associations by treating them as candidate predictive features rather than endpoints, and it differs by predicting full FLAME shape. The work helps FaceCon justify a GWAS-informed panel and provides the comparison point for the ablation against matched-random SNPs.

\subsubsection{Naqvi et al. --- Shared Heritability of Face and Brain Shape}
Summary. Naqvi et al.~\cite{ref:naqvi} investigate the genetic overlap between facial and brain shape. Using genome-wide analysis and genetic-correlation methods on European cohorts, the study shows that many loci influencing the face also relate to cranial and neuroanatomical variation, deepening the understanding of the developmental genetics that couples facial form to underlying structure.

Critical Analysis. The work's strength is that it establishes the biological embedding of facial genetics within broader craniofacial development, which supports the plausibility of predicting shape from genotype. Its relevance to prediction is indirect, however, because it characterizes genetic overlap rather than reconstructing faces, and like the other studies it is concentrated in a single ancestry group.

Relationship to the Proposed Work. This study provides part of the biological rationale that FaceCon cites for attempting genotype to shape prediction, and it informs which loci are most likely to carry shape-relevant signal. FaceCon does not adopt a method from it, but uses it to motivate the scientific premise and to reinforce the expectation that the predictable signal is diffuse.

\subsubsection{Xiong et al. --- Combined GWAS and Facial Prediction}
Summary. Xiong et al.~\cite{ref:xiong} present a combined genome-wide association study of facial traits in Europeans whose explicit objective is to increase explained variance and to improve facial prediction. The methodology analyses close to one thousand facial measurements in over eleven thousand individuals, combining cohorts to boost power, and couples association discovery with prediction modelling. The study reports a large number of associated genomic regions, including many previously unknown, explains substantially more genetic variance than earlier work, and demonstrates improved facial prediction, extending even to the reconstruction of archaic human facial profiles.

Critical Analysis. This is among the most prediction-oriented facial-genetics studies, and its explicit treatment of prediction and cross-population validation is a model for careful reporting. Its limitations are that it operates on facial measurements and landmarks rather than full generative 3D meshes, and that it remains concentrated in European ancestry, so its prediction performance may not transfer directly to other populations.

Relationship to the Proposed Work. Xiong et al. serve FaceCon both as an SNP-selection resource and as a methodological benchmark for honest, prediction-focused reporting. FaceCon extends this line of work from measurement-level prediction into generative 3D FLAME prediction, and it adopts the study's discipline of reporting prediction against clear reference points, which FaceCon implements as its demographic baselines.

\subsubsection{Richmond et al. --- Overview of Facial Genetics}
Summary. Richmond et al.~\cite{ref:richmond} provide an accessible review of facial genetics. The work summarizes heritability evidence, the growth of facial genome-wide association studies, and the clinical and forensic relevance of predicting appearance from DNA, consolidating the field for a broad audience.

Critical Analysis. As a review, its strength is a balanced, consolidated overview that situates the individual studies within a coherent picture. Its limitation is that it contributes no new model or dataset, so it informs context rather than method.

Relationship to the Proposed Work. FaceCon uses this review to frame the problem domain and to justify the relevance of genotype to face research. It adopts no technique from it, but draws on it to motivate the project and to set realistic expectations about what facial genetics can and cannot currently achieve.

Synthesis. Collectively, these studies establish that facial morphology is genetically influenced and polygenic, and they provide reproducible, published associations between variants and facial traits. They share two limitations that shape FaceCon directly: they are largely discovery- or measurement-oriented rather than generative, and they are ancestry-concentrated. None provides an open, generative 3D predictor. This is the research gap that FaceCon targets, and it is why FaceCon builds a transparent, GWAS-informed SNP panel and commits in advance to ancestry-stratified evaluation.

\subsection{Forensic DNA Phenotyping}
This topic covers the applied setting in which appearance is predicted from DNA, which is where genotype to face research is ultimately motivated.

\subsubsection{Wang et al. --- A 19-SNP Forensic Prediction System}
Summary. Wang et al.~\cite{ref:wang} describe a forensic DNA-phenotyping system based on a compact nineteen-SNP panel for predicting externally visible characteristics. The work follows the established forensic pattern of selecting and validating a small, interpretable panel for categorical appearance traits so that predictions can support investigations from unidentified genetic material.

Critical Analysis. The strength of this approach is its rigor and deployability: a small validated panel is interpretable, inexpensive to genotype, and suited to casework. Its limitation, from FaceCon's perspective, is that it predicts discrete traits such as pigmentation rather than continuous three-dimensional facial geometry, so it addresses a much narrower output than full facial shape.

Relationship to the Proposed Work. FaceCon shares the forensic and anthropological motivation of this work and targets the harder output of continuous 3D facial shape, while integrating, rather than retraining, established forensic phenotyping to provide the categorical colour traits (eye, hair, and skin colour) alongside the reconstructed face. It adopts the forensic discipline of small, documented, validated panels for appearance while remaining explicitly scoped as research with firm limits on individual-level claims. The work helps define both the appearance method FaceCon integrates and the shape gap it explores.

\subsubsection{HIrisPlex-S --- Eye, Hair, and Skin Colour Prediction}
Summary. The HIrisPlex and HIrisPlex-S systems~\cite{ref:hirisplex,ref:hirisplexs} predict eye and hair colour, and subsequently skin colour, from compact validated SNP panels, and are the de facto standard for appearance prediction in forensic genetics. Complementary genome-wide association work in admixed cohorts~\cite{ref:adhikari} further characterises the genetic basis of human pigmentation across diverse populations.

Critical Analysis. These systems are rigorously validated and report calibrated, honest performance, which is strong at the pigmentation extremes (for example blue and brown eyes, pale and dark skin, red and black hair) and genuinely weaker for intermediate categories. Their scope is categorical appearance rather than 3D geometry, and performance is partly entangled with ancestry.

Relationship to the Proposed Work. FaceCon's appearance component follows this literature by integrating the established HIrisPlex-S predictor directly: the same genotype supplies HIrisPlex-S's validated marker SNPs, and the external predictor returns eye colour, hair colour, and skin colour, which are shown alongside the reconstructed face. FaceCon does not retrain these forensic models or extend them to other traits, and it reports their colours using the published method rather than re-deriving accuracy on its own data.

Synthesis. The forensic literature demonstrates both the demand for genotype to appearance prediction and the discipline with which it is usually approached, through small validated panels, categorical traits, and honest calibration. FaceCon incorporates this appearance line of work as a dedicated module and combines it with continuous 3D shape prediction, positioning itself across both while refusing to present itself as an operational identification tool.

\subsection{Statistical and Learning Methods for Genotype to Face Prediction}
This topic covers approaches that connect genotype and face using statistical and learning models, which form FaceCon's essential baselines and its closest methodological neighbours.

\subsubsection{Sero et al. --- Face to DNA Classifiers}
Summary. Sero et al.~\cite{ref:sero} take a complementary, face to DNA direction. The objective is to quantify how much genomic information is recoverable from a face by building classifiers that predict DNA-encoded attributes (such as genomic background and individual loci) and DNA-inferred attributes (such as age) from three-dimensional faces, and then matching individuals between the two modalities. The data pair 3D facial images with genomic data across both diverse and more homogeneous cohorts, and the study reports verification-mode matching rates well above chance, on the order of eighty percent.

Critical Analysis. The study is rigorous and quantitative, and its most valuable feature is its honesty: the authors show that much of the matching signal is carried by broad attributes such as gender and ancestry rather than by fine individual genetics. Its limitations are that it works in the reverse (face to DNA) direction rather than generating a face, that verification is weaker than individual identification, and that the demographic contribution it exposes is a warning about how easily apparent performance can be overstated.

Relationship to the Proposed Work. This work is, in effect, a specification for part of FaceCon's evaluation. FaceCon adopts its central lesson by making demographic baselines mandatory and by decomposing performance into demographic and individual-genetic components. FaceCon differs by predicting in the forward direction, from genotype to 3D shape, and by using the demographic decomposition not as an afterthought but as the headline measure of success.

Synthesis. Statistical and learning methods are interpretable, sample-efficient, and often competitive when paired sample sizes are modest, which is exactly the regime FaceCon expects. The decisive lesson of this topic is that demographic attributes can dominate apparent performance unless they are explicitly controlled, which is why FaceCon adopts ridge regression and demographic regressors as mandatory baselines and builds the demographic-versus-genetic decomposition directly into its evaluation.

\subsection{Deep Learning for Genotype to Face Generation}
This topic covers the neural systems that treat genotype to face as multimodal generation, together with the critique and the broader guidance that shape how such systems should be built and judged.

\subsubsection{Jiao et al. --- Difface}
Summary. Jiao et al.~\cite{ref:difface} propose Difface, the most visible recent system for DNA to 3D face generation. The objective is to reconstruct a dense three-dimensional facial mesh from DNA alone. Methodologically, Difface is a two-stage, multi-modality model: it first contrastively aligns a transformer SNP encoder with a mesh face encoder, and it then uses a diffusion prior in the shared embedding space to reconstruct the mesh. It was trained on roughly 9,674 Han Chinese subjects using about 7,842 GWAS-selected SNPs, and it demonstrates that modern generative models can map genotype to a recognizable 3D face.

Critical Analysis. Difface is a compelling proof of concept and combines several strong ideas, but its limitations are decisive for reproducibility and for scientific interpretation. Its public release does not provide the SNP list, the full preprocessing, the face template, or reusable pretrained weights, so it cannot be independently reproduced. Its evaluation does not cleanly separate individual-level prediction from demographic averages, its cohort is drawn from a single ancestry group, and its dense-mesh target is high-dimensional and tied to an undocumented topology.

Relationship to the Proposed Work. FaceCon treats Difface as an architectural reference rather than a system to reproduce. It adopts the high-level two-stage alignment and generation idea, but it re-implements the pipeline on documented 3DFN and FaceBase data, replaces the undisclosed panel with a documented SNP-selection procedure, substitutes the shared FLAME representation for the unknown dense mesh, and, most importantly, adds the demographic baselines whose absence is Difface's central evaluation weakness. FaceCon therefore directly targets the reproducibility and confounding gaps that Difface leaves open.

\subsubsection{Wagner et al. --- Critique of Difface}
Summary. Wagner et al.~\cite{ref:comment} respond to Difface with a critical commentary. The work argues that facial-genetics AI systems require transparent metrics, clear disclosure of data flow, and evaluation that explicitly separates demographic signal from individual genetic prediction, and it uses Difface as a case study for these requirements.

Critical Analysis. The commentary's strength is that it names, precisely, the methodological trap that the whole subfield must avoid, turning a general concern into a concrete evaluation requirement. Its limitation is that, as a commentary, it diagnoses the problem rather than providing an open system that demonstrates the solution.

Relationship to the Proposed Work. FaceCon operationalizes this critique. The requirements it raises, transparency, disclosed data flow, and confounder separated evaluation, become concrete design commitments in FaceCon's methodology, and in particular its demographic-baseline-centered evaluation is a direct response to the concerns raised here.

\subsubsection{Alharbi and Rashid --- Deep Learning in Human Genomics}
Summary. Alharbi and Rashid~\cite{ref:alharbi} survey deep learning in human genomics using next-generation sequencing data. The review maps the architectures, tasks, and recurring challenges of the field, with particular attention to convolutional, recurrent, and attention models and to the practical problems of data scarcity, confounding, and interpretability.

Critical Analysis. The survey's strength is an up-to-date and balanced account of methods and pitfalls, which is directly relevant to a small-sample genotype to face project. Its limitation is that it is broad and does not address the specific genotype to 3D shape task, so its lessons must be adapted rather than applied directly.

Relationship to the Proposed Work. This review justifies FaceCon's methodological caution. Its warnings about data scarcity and confounding support FaceCon's decision to prefer compact models, to evaluate complex architectures only against simpler baselines, and to treat interpretability and honest reporting as first-class concerns.

Synthesis. FaceCon treats Difface as an architectural reference, particularly its two-stage alignment and generation framework, while adapting the modelling approach to the 3DFN dataset and FLAME representation. It incorporates the evaluation considerations raised by Wagner et al. and the cautions of Alharbi and Rashid regarding model complexity and limited sample sizes.

\subsection{3D Facial Representations and FLAME}
This topic covers how a 3D face can be represented so that it is comparable across subjects and usable as a learning target.

\subsubsection{Egger et al. --- Survey of 3D Morphable Models}
Summary. Egger et al.~\cite{ref:egger} survey two decades of 3D morphable face models, covering how they are constructed, fitted, and applied, and situating modern parametric models within that history. The survey provides a comprehensive and authoritative map of the representational landscape in which FLAME sits.

Critical Analysis. Its strength is breadth and authority, giving a balanced account of what morphable models can and cannot do. Its limitation is that, being published before the most recent neural generative advances, it only anticipates some of the directions that are now central.

Relationship to the Proposed Work. This survey frames FaceCon's representational choice and clarifies the trade-offs, particularly the loss of fine detail, that come with a compact parametric target. It justifies, at a conceptual level, FaceCon's decision to predict a morphable-model parameterization rather than raw geometry.

\subsubsection{Li et al. --- The FLAME Head Model}
Summary. Li et al.~\cite{ref:flame} propose FLAME, a parametric head model with a linear identity shape space, expression blendshapes, and pose, learned from thousands of 3D head scans. FLAME represents any face as a short parameter vector and decodes that vector back into a full mesh with a fixed, shared topology.

Critical Analysis. FLAME's strengths are topological consistency and a compact, interpretable parameterization that is well suited to learning. Its limitations are a ceiling on fine detail inherent to a linear shape space and a dependence on the demographic composition of its training scans, which can bias the representable range of faces.

Relationship to the Proposed Work. FLAME is the backbone of FaceCon's face representation. FaceCon adopts FLAME directly: it registers every FaceBase scan to FLAME so that all subjects share a topology, and it predicts FLAME identity parameters as its learning target. This choice is what makes the generative problem compact and comparable across models.

\subsubsection{Feng et al. --- DECA}
Summary. Feng et al.~\cite{ref:deca} present DECA, which learns to regress a detailed, animatable 3D face model from unconstrained single images. Building on a FLAME-style parameterization, DECA uses a detail-consistency loss to separate person-specific shape detail from expression-dependent variation, and it recovers convincing geometry even from weak, in-the-wild inputs.

Critical Analysis. DECA's strength is its demonstration that a FLAME target pairs effectively with a learned encoder under difficult input conditions. Its limitation, for FaceCon, is that its input modality is images rather than genotype, and its fine detail relies on image-texture cues that have no analogue in DNA, so only its representational strategy transfers.

Relationship to the Proposed Work. DECA validates the design choice at the heart of FaceCon's output, namely pairing a FLAME parameterization with a learned encoder-decoder. Where DECA regresses FLAME from pixels, FaceCon generates FLAME from genotype, but both rely on the same compact, topology consistent target, and DECA's success on weak inputs is encouraging for an even weaker conditioning signal.

\subsubsection{Gong et al. --- SpiralNet++}
Summary. Gong et al.~\cite{ref:spiralnet} propose SpiralNet++, an efficient spiral mesh-convolution operator for neural networks that consume meshes of fixed topology. It defines local neighbourhoods along spirals on a mesh and aggregates them efficiently, providing a simple and fast operator for mesh learning.

Critical Analysis. Its strength is efficiency and simplicity for fixed-topology meshes. Its limitation is precisely that it assumes a fixed vertex ordering, which is only available after registration to a common template, so it cannot be applied to raw, uncorresponded scans.

Relationship to the Proposed Work. SpiralNet++ becomes applicable to FaceCon only because every scan is registered to the shared FLAME topology. It is a candidate operator for an advanced mesh encoder, and it illustrates why the FLAME registration step is a prerequisite for mesh-convolutional models within FaceCon.

Synthesis. This topic supplies FaceCon's output representation. Morphable models, and FLAME in particular, give the topological consistency and compactness that make 3D faces learnable; DECA shows the target pairs well with learned encoders; and SpiralNet++ becomes usable once scans are registered to a common FLAME topology. FaceCon therefore fits every FaceBase scan to FLAME and predicts FLAME shape parameters, accepting reduced fine detail in exchange for a tractable, comparable learning target.

\subsection{Advanced Representation and Generative Modeling Methods}
For the advanced modelling stage, FaceCon considers modern representation-learning and generative approaches. Denoising diffusion probabilistic models~\cite{ref:ddpm} learn to reverse a gradual noising process to generate samples, while latent diffusion~\cite{ref:ldm} performs this process in a learned lower-dimensional representation to reduce computational cost and support conditional generation. Contrastive learning, as popularized by CLIP~\cite{ref:clip}, can align two different modalities within a shared feature space. Transformer architectures~\cite{ref:transformer} use self-attention to model interactions across input features and can be adapted to SNP data using suitable embeddings and genomic positional information.

Critical Analysis. These methods provide greater modelling capacity but also introduce higher data and computational requirements. This is particularly relevant for genotype to face prediction, where the number of paired subjects is small relative to the dimensionality of the genetic input. More complex architectures must therefore be evaluated against simpler statistical and neural baselines to determine whether their additional capacity produces meaningful improvements.

Relationship to the Proposed Work. FaceCon follows a progressive modelling strategy from statistical SNP models to MLPs and transformer models. Contrastive alignment will be investigated for learning shared SNP-face representations, while generative approaches such as diffusion will be evaluated where the size of the final paired dataset and preceding experimental results support their use.

\subsection{Datasets and Tooling}
This topic covers the data resources and tools that make a reproducible pipeline possible. The 3D Facial Norms Database provides high-quality craniofacial anthropometry and 3D facial surface models through the FaceBase consortium~\cite{ref:facebase}, and genotypes for related normal facial variation are available through dbGaP accession phs000949~\cite{ref:dbgap}. Genotype quality control and SNP handling are supported by standard, community-maintained tooling~\cite{ref:plink}. The strength of this setting is that it is documented and reusable, unlike the private cohorts of closed systems, while its limitations are modest sample size and ancestry concentration. These resources are FaceCon's data foundation, and using them is precisely what allows the whole pipeline to be reproduced by others.

\section{Literature Review Summary Table}
Table~\ref{tab:litsummary} summarizes the reviewed research in the research-project format of author, method, results or contribution, and limitations.

{\small\setlength{\LTpre}{6pt}\setlength{\LTpost}{0pt}
\begin{longtable}{|p{2.6cm}|p{3.5cm}|p{4.0cm}|p{3.2cm}|}
\caption{Summary of the literature reviewed for FaceCon (research-project format).}\label{tab:litsummary}\\
\hline
\textbf{Author} & \textbf{Method} & \textbf{Results / Contribution} & \textbf{Limitations} \\ \hline
\endfirsthead
\hline
\textbf{Author} & \textbf{Method} & \textbf{Results / Contribution} & \textbf{Limitations} \\ \hline
\endhead
\hline
\endfoot
Claes et al.~\cite{ref:claes} & Global to local segmented GWAS & Many novel facial-shape loci; localized effects & Discovery-oriented; ancestry-concentrated \\ \hline
White et al.~\cite{ref:white} & Multivariate facial GWAS & Expanded catalogue of face-associated variants & Mostly European cohorts \\ \hline
Naqvi et al.~\cite{ref:naqvi} & GWAS / shared-heritability analysis & Links face and brain shape genetics & Indirect for prediction \\ \hline
Xiong et al.~\cite{ref:xiong} & Combined GWAS + prediction & More explained variance; improved prediction & Landmark/measurement; European \\ \hline
Richmond et al.~\cite{ref:richmond} & Review of facial genetics & Accessible overview of loci and relevance & Not a predictive model \\ \hline
Wang et al.~\cite{ref:wang} & 19-SNP forensic prediction system & Validated categorical-trait prediction & Categorical traits only; not 3D shape \\ \hline
Sero et al.~\cite{ref:sero} & Face to DNA classifiers & Above-chance verification matching & Signal partly demographic; not full 3D generation \\ \hline
Jiao et al.~\cite{ref:difface} & Transformer + mesh encoder; diffusion prior & First de novo 3D-face reconstruction from DNA & Private data; undisclosed panel/weights; weak demographic control \\ \hline
Wagner et al.~\cite{ref:comment} & Critical commentary & Specifies transparent, confounder aware evaluation & Diagnoses rather than provides a system \\ \hline
Alharbi \& Rashid~\cite{ref:alharbi} & Survey of deep learning in genomics & Maps methods, pitfalls, confounding & Not specific to 3D-shape prediction \\ \hline
Egger et al.~\cite{ref:egger} & Survey of 3D morphable models & Comprehensive map of the field & Pre-dates recent neural methods \\ \hline
Li et al.~\cite{ref:flame} & FLAME parametric head model & Compact, topology consistent shape space & Linear; limited fine detail; population bias \\ \hline
Feng et al.~\cite{ref:deca} & FLAME + learned encoder (DECA) & Detailed 3D face from in-the-wild images & Image inputs, not genotype \\ \hline
Gong et al.~\cite{ref:spiralnet} & SpiralNet++ mesh convolution & Fast, efficient mesh operator & Requires fixed topology \\ \hline
Rombach et al.~\cite{ref:ldm} & Latent diffusion & Efficient, controllable generation & Needs large data; conditioning quality matters \\ \hline
Radford et al.~\cite{ref:clip} & Contrastive multimodal alignment & Shared cross-modal embedding space & Scales with paired-data volume \\ \hline
Ho et al.~\cite{ref:ddpm} & Denoising diffusion (DDPM) & High-quality generative modeling & Data- and compute-intensive \\ \hline
Vaswani et al.~\cite{ref:transformer} & Transformer self-attention & Strong sequence/feature encoder & Data-hungry on small panels \\ \hline
Chang et al.~\cite{ref:plink} & Second-generation PLINK & Standard genotype QC toolkit & Scientific thresholds left to user \\ \hline
\end{longtable}
}

\section{Conclusion}
Taken together, the literature both justifies FaceCon and constrains its design. The facial-GWAS studies establish that facial shape is genetically influenced and supply a transparent evidence base for SNP selection, yet they are discovery- or measurement-oriented and ancestry-concentrated rather than generative. Forensic phenotyping and face to DNA work show the real demand for genotype to appearance prediction, but they also expose the central danger of the field, which is that demographic attributes such as age, gender, and ancestry can masquerade as individual genetic signal. Difface demonstrates that end-to-end DNA to 3D face generation is feasible, while its published critique and the broader genomics deep learning survey make clear that such systems must be reproducible and must evaluate against demographic baselines. The morphable-model literature supplies the shared FLAME representation that makes the learning target tractable, and the advanced generative methods offer additional capacity that is only justified once simpler models have shown that learnable signal exists.

The research gap is therefore specific: there is no open, reproducible, and demographically controlled baseline that predicts 3D facial shape from genotype on documented data. This gap directly justifies FaceCon's approach. It justifies building the pipeline on documented 3DFN and FaceBase data, using FLAME as the common representation, and, above all, it justifies the progression from simpler to more complex models. Because the data are small, high-dimensional, and confounded, complex architectures cannot be trusted until they have been shown to beat simple and demographic baselines on identical splits and metrics, which is exactly why FaceCon evaluates simpler baselines before more complex models. The next chapter translates these commitments into a software requirements specification.
